\section{Function information, derivative constructions and certified roots}
\label{sec:p04-computable-real-analysis}

The distinction between a value oracle and a uniform function name is
essential in numerical reasoning. The representations and all-name
realizer convention are those of
\cref{sec:p04-represented-spaces,sec:p04-function-operators}.
In this section \(\delta_{\mathrm C}\) still names scalar reals by
pointwise-error Cauchy approximations, and \(\gamma_{\mathrm P}\)
names members of \(C[0,1]\) by rational polygons with uniform error
at most \(2^{-n}\). Neither representation is replaced by an
uncontrolled convergent sequence.
The modulus and operator constructions extend the explicitly represented
procedures in \cite[Theorems 6.4, 6.7, 6.9 and 7.7]{Weihrauch1995IntroductionRS}.
The differentiability construction is a corrected reconstruction of the bump presentation in
\cite[\S8.2, pp.\ 50--53]{Bournez2024ComputableAnalysisAuthor},
with the explicit parameters and full error estimates below.

\subsection{Rational samples and the information supplied by a modulus}
\label{sec:p04-ca-samples}

\begin{hbdefinition}{def:p04-ca-sample-name}{A sampled-value representation}
Fix an effective enumeration \((q_i)_{i\in\mathbb N}\) of all
rationals in \([0,1]\), and an effective enumeration of pairs
\((i,n)\in\mathbb N^2\). A \(\sigma_{\mathrm S}\)-name of
\(f\in C[0,1]\) gives a rational \(r_{i,n}\) at the position
for \((i,n)\), with
\[
 |r_{i,n}-f(q_i)|\leq2^{-n}.
\]
The entries use the rational coding of \cref{def:p04-rs-rational-code}.
The continuity promise makes the named function unique, since two
continuous functions agreeing on all rational arguments agree everywhere.
A modulus is a function \(\mu:\mathbb N\to\mathbb N\) with
\[
 |x-y|\leq2^{-\mu(n)}\ \Longrightarrow\
 |f(x)-f(y)|\leq2^{-n}\qquad(x,y\in[0,1]).
\]
A sample-and-modulus name consists of the sample table and a stream
giving the integer values of such a modulus. The promise applies to
the same function in both components.
\end{hbdefinition}

\begin{hbtheorem}{thm:p04-ca-sample-modulus}{proved=annex:p04-ca-sample-modulus}{Uniform names and sample tables with a supplied modulus}
Uniform polygon names are computably equivalent to sample-and-modulus
names. In the direction from samples to a polygon of output precision
\(n\), use a dyadic partition of mesh
\(2^{-\mu(n+2)}\), approximate the values at its vertices within
\(2^{-n-2}\), and interpolate those rational values linearly.
The resulting uniform error is at most \(2^{-n-1}\).
The converse computes the table and a modulus from the polygon name.
\end{hbtheorem}

This equivalence is between kinds of input information. Merely knowing
that a continuous function has a modulus does not supply the integer
values of a computable modulus or compute them from a weak name.

\begin{hbtheorem}{thm:p04-ca-sample-obstruction}{proved=annex:p04-ca-sample-obstruction}{Samples alone do not uniformly determine area or a uniform name}
On all of \(C[0,1]\), the scalar operation
\(f\mapsto\int_0^1 f(x)\,dx\) has no continuous
\((\sigma_{\mathrm S},\delta_{\mathrm C})\)-realizer.
Consequently there is no continuous translator from sampled-value names
to uniform polygon names, and no computable translator of that kind.
The assertion still fails when each queried rational value is supplied
exactly for a rational polygon: finitely many samples can miss a narrow,
high tent with area one.
\end{hbtheorem}

The positive integration result in \cref{thm:p04-fo-integration} is not
in conflict with this obstruction: its input controls the whole function
uniformly. A supplied height bound together with an appropriate modulus
changes the available information; neither is silently included in a
sample table.

\subsection{A computable function with a noncomputable continuous derivative}
\label{sec:p04-ca-bumps}

\begin{hbdefinition}{def:p04-ca-bump}{A differentiable bump and its scales}
Define, on \(\mathbb R\),
\[
 \theta(x)=
 \begin{cases}x(x^2-1)^2,&-1\leq x\leq1,\\
 0,&|x|>1,\end{cases}
 \qquad \lambda=\frac{16}{25\sqrt5},\qquad
 \theta_{\alpha\beta}(x)=\frac{\beta}{\lambda}
                         \theta(x/\alpha),
\]
where \(\alpha,\beta>0\). Here \(\lambda\) is the maximum
of \(|\theta|\); \(\alpha\) is a half-width and \(\beta\)
an amplitude, not derivative bounds. The basic bump is odd,
vanishes at \(-1,0,1\), and has zero derivative at \(-1,1\)
and derivative one at zero. Its minimum and maximum occur at
\(-1/\sqrt5\) and \(1/\sqrt5\). In particular,
\[
 \lVert\theta_{\alpha\beta}\rVert_\infty=\beta,
 \qquad \theta'_{\alpha\beta}(0)=\frac{\beta}{\lambda\alpha},
 \qquad |\theta'_{\alpha\beta}(x)|\leq
                         \frac{\beta}{\lambda\alpha}.
\]
The derivative at each support endpoint is zero.
\end{hbdefinition}

The function is computable without deciding whether an arbitrary real
argument equals a support endpoint. Indeed, with
\(c(x)=\min(1,\max(-1,x))\), one has
\(\theta(x)=c(x)(c(x)^2-1)^2\). Clamping, polynomial arithmetic,
and the positive computable constant \(\lambda\) give a uniform
value algorithm.

\begin{hbdefinition}{def:p04-ca-coding-function}{Widths, heights and an injective enumeration}
Let \(A\subseteq\mathbb N\) be computably enumerable but undecidable.
Fix a total computable injective enumeration
\(h:\mathbb N\to\mathbb N\) with range \(A\).
Such a set is infinite, and waiting for the next previously unseen member
of an enumeration constructs \(h\).
For \(n=h(k)\), define
\[
 \alpha_n=2^{-k-2n-2},\qquad
 \beta_n=2^{-k-3n-2},\qquad
 \frac{\beta_n}{\alpha_n}=2^{-n}.
\]
These symbols are defined for \(n\in A\); finding its enumeration
position is not an operation promised for an arbitrary \(n\).
Define, for \(x\in[0,1]\),
\[
 f(x)=\sum_{k=0}^{\infty}
   \theta_{\alpha_{h(k)}\beta_{h(k)}}(x-2^{-h(k)}),
 \qquad
 g_M(x)=\sum_{k=0}^{M}
   \theta_{\alpha_{h(k)}\beta_{h(k)}}(x-2^{-h(k)}).
\]
The supports are disjoint, so at each point at most one summand is
nonzero. The notation is a series, but its pointwise value is not an
uncontrolled infinite accumulation of nonzero terms.
\end{hbdefinition}

\begin{hbtheorem}{thm:p04-ca-computable-c1}{proved=annex:p04-ca-computable-c1}{A full uniform name for the bump construction}
The function just defined is computable in the uniform polygon
representation and belongs to \(C^1[0,1]\).
Its derivative is continuous at the accumulating endpoint, with
\(f'(0)=0\). The function \(f\) is Lipschitz with constant four.
Its second right derivative at zero does not exist.
The finite partial sums satisfy
\[
 \lVert f-g_M\rVert_\infty\leq2^{-M-3}.
\]
For every \(n\),
\[
 f'(2^{-n})=
 \begin{cases}2^{-n}/\lambda,&n\in A,\\0,&n\notin A.
 \end{cases}
\]
Endpoint derivatives are one-sided on the closed interval.
\end{hbtheorem}

\begin{hbtheorem}{thm:p04-ca-noncomputable-derivative}{proved=annex:p04-ca-noncomputable-derivative}{One function, not only a noncomputable operator}
The continuous derivative of this computable \(C^1\) function has no
computable uniform polygon name. It has no computable modulus of
continuity. Computing the derivative uniformly at the rational centres
\(2^{-n}\) would decide \(A\), using the separated values in the
preceding theorem.
\end{hbtheorem}

This is stronger than the uniform differentiation obstruction in
\cref{thm:p04-fo-differentiation}: it exhibits one computable input whose
derivative is not a computable function. Each specified centre value is
nevertheless individually a computable scalar, either zero or
\(2^{-n}/\lambda\). There is no uniform procedure selecting the correct
alternative from \(n\). Function computability is not the assertion that
each value separately has some program.
In contrast, \(f(2^{-n})=0\) in both cases. The nearby function heights
depend on the enumeration time \(k\) and can be arbitrarily small;
the derivative separation does not imply a corresponding decidable
test on the original function's values.

\subsection{Additional differentiability and additional input information}
\label{sec:p04-ca-higher-derivatives}

\begin{hbtheorem}{thm:p04-ca-derivative-modulus}{proved=annex:p04-ca-derivative-modulus}{The modulus criterion and its individual \(C^2\) consequence}
For a fixed computable \(f\in C^1[0,1]\), its derivative is
computable if and only if its derivative has a computable modulus.
For a fixed computable \(f\in C^2[0,1]\), the derivative is therefore
computable. More explicitly, if a positive integer bound \(K\) on
\(\lVert f''\rVert_\infty\) is supplied, choose \(\ell\) with
\(K\leq2^\ell\); then \(\mu'(n)=n+\ell\) is a computable
derivative modulus and computes \(f'\) uniformly from \(f\)'s name
and that bound. Every derivative of one fixed computable \(C^\infty\)
function is individually computable.
\end{hbtheorem}

The first clause completes the exercise stated as Theorem 8.6 in
\cite[pp.\ 50--51]{Bournez2024ComputableAnalysisAuthor}; the remaining
clauses make the information in its Corollaries 8.1--8.2 explicit.
Compactness supplies existence of a bound on a continuous second
derivative, not a procedure that computes that bound from an ordinary
name of \(f\). The individual \(C^2\) conclusion permits the existing
integer bound to be fixed in a program for this one function.
It does not supply a uniform realizer on all ordinary names of
\(C^2\) functions. Likewise, separate computability for each derivative
order does not supply one program that accepts the order as input and
produces all the derivative names.

\subsection{Turning a small residual into a certified root location}
\label{sec:p04-ca-root-certificates}

\begin{hbdefinition}{def:p04-ca-root-bound}{A lower-slope promise}
Let \(f\in C^1[0,1]\), let \(m>0\) be supplied rational data,
and promise \(f'(x)\geq m\) on \([0,1]\) and
\(f(0)\leq0\leq f(1)\). There is then exactly one zero \(z\).
The lower-slope condition is not established by observing finitely
many samples, nor by the bare fact that the zero is unique.
\end{hbdefinition}

\begin{hbtheorem}{thm:p04-ca-root-certificate}{proved=annex:p04-ca-root-certificate}{Residual-to-location and perturbation bounds}
Under the lower-slope promise, every \(x\in[0,1]\) satisfies
\[
 |x-z|\leq |f(x)|/m.
\]
From a uniform polygon name \((P_k)\) and output precision \(n\),
search for rational \(s\in[0,1]\) and \(k\) with
\[
 |P_k(s)|+2^{-k}<m2^{-n}.
\]
The search terminates, and \(s\) is a root approximation with
\(|s-z|<2^{-n}\). If a second continuous function \(g\) has
a zero \(w\in[0,1]\) and
\(\lVert f-g\rVert_\infty\leq\epsilon\), then
\(|w-z|\leq\epsilon/m\).
\end{hbtheorem}

For \(f(x)=x-1/3\), take \(m=1\): the residual is exactly
the location error. For \(f(x)=\eta(x-1/3)\), with small
positive \(\eta\), the same location error has residual smaller by
the factor \(\eta\), and the valid certificate uses \(m=\eta\).
For \(f(x)=(x-1/2)^2\), uniqueness holds but no positive lower-slope
bound does: near the root, distance is the square root of the residual,
not the residual divided by a fixed positive constant. The unique-root
algorithm in \cref{thm:p04-fo-unique-zero} still applies; its correctness
does not depend on a nonexistent slope certificate.
